% Diagramas vetoriais: estilos comuns, legíveis também em escala de cinzentos.
\usepackage{tikz}
\usetikzlibrary{arrows.meta,positioning,calc,fit,backgrounds}
\definecolor{LasetBlue}{HTML}{234D65}
\definecolor{LasetTeal}{HTML}{287D83}
\definecolor{LasetTint}{HTML}{EDF3F6}
\definecolor{LasetGrey}{HTML}{555E64}
\tikzset{
    bloco/.style={draw=LasetBlue,fill=LasetTint,rounded corners=2pt,
        line width=0.7pt,align=center,inner sep=5pt,
        font=\sffamily\small,text=black},
    nucleo/.style={bloco,fill=LasetBlue,text=white},
    auxiliar/.style={bloco,draw=LasetGrey,fill=black!3},
    opcional/.style={auxiliar,dashed},
    fluxo/.style={-{Latex[length=2.2mm,width=1.5mm]},
        draw=LasetBlue,line width=0.85pt},
    duplo/.style={{Latex[length=2mm,width=1.4mm]}-{Latex[length=2mm,width=1.4mm]},
        draw=LasetBlue,line width=0.85pt},
    ligacao/.style={draw=LasetBlue,line width=0.8pt},
    legenda/.style={font=\sffamily\footnotesize,align=center,text=LasetGrey},
    mini/.style={bloco,font=\sffamily\footnotesize,inner sep=3pt}
}
\newcommand{\Diagrama}[1]{%
    \begingroup\setstretch{1}\hyphenpenalty=10000\exhyphenpenalty=10000%
    \input{./Figuras/diagramas/#1.tex}\endgroup}
